% Source: https://artofproblemsolving.com/community/c3h1388540_we_all_scream_for_ice_cream
% Stable identity: aops-topic-1388540-q01
\textbf{The frustum.} Extend the slanted sides of the cup down to the apex $V$ of the full cone (dashed). The radii $1$ and $2$ are in ratio $1:2$, so the small cone cut off below the cup has height $4$ and the full cone has height $8$.

\begin{asy}
usepackage("amsmath");
size(4.5cm);
defaultpen(fontsize(10pt));
real k=0.28;
pair T=(0,8);
draw((-2,4)--(-1,0),linewidth(0.9)); draw((2,4)--(1,0),linewidth(0.9));
draw((-2,4)--T--(2,4),linewidth(0.9));
draw(yscale(k)*arc((0,0),1,180,360)); draw(yscale(k)*arc((0,0),1,0,180),dashed);
draw(shift(0,4)*yscale(k)*scale(2)*unitcircle);
draw((-1,0)--(0,-4)--(1,0),gray+dashed);
draw((2.6,0)--(2.6,4),Arrows); label("$4$",(2.6,2),E);
draw((2.6,4)--(2.6,8),Arrows); label("$4$",(2.6,6),E);
draw((0,0)--(-1,0)); label("$1$",(-1,0),W,fontsize(9pt));
draw((0,4)--(-2,4)); label("$2$",(-2,4),W,fontsize(9pt));
dot((0,-4)); label("$V$",(0,-4),S);
label("\mbox{}",(min(currentpicture,user=true)+max(currentpicture,user=true))/2,black); // keeps the text after the figure black when LaTeX typesets the labels inline
\end{asy}

So the frustum's volume is
\[
\dfrac{1}{3}\pi\cdot2^2\cdot8-\dfrac{1}{3}\pi\cdot1^2\cdot4=\dfrac{32\pi}{3}-\dfrac{4\pi}{3}=\dfrac{28\pi}{3}.
\]
(The frustum formula $\dfrac{\pi h}{3}\left(R^2+Rr+r^2\right)=\dfrac{4\pi}{3}(4+2+1)$ gives the same.)

\textbf{The cone on top} has radius $2$ and height $4$: volume $\dfrac{1}{3}\pi\cdot2^2\cdot4=\dfrac{16\pi}{3}$.

The total volume is $\dfrac{28\pi}{3}+\dfrac{16\pi}{3}=\boxed{\dfrac{44\pi}{3}}$ cubic inches.
